
\chapter{Repairing the graph memory pipeline: what improved and what remained}
\label{ch:graph_mechanisms}

On 200 LoCoMo questions, the initial Gemma~4 E2B graph baseline scored 15.0\%, against
36.0\% for vector retrieval.\footnote{The earlier unprefixed vector-retrieval runs scored
32.5\%.} We investigated whether
the graph deficit came from implementation problems
or from the evidence supplied by the graph. A repair sequence changed retrieval, increased
the extraction budget and restored source dialogue. Further controls varied graph depth
and separated the extractor from the reader.

The repairs improved accuracy, with the largest single gain following source restoration.
This chapter examines where evidence was lost and which repairs recovered it, motivating
the budgeted source-selection policy evaluated in Chapter~\ref{ch:spe}.

\paragraph{Contribution of this chapter.} \citet{rag_vs_graphrag} compare triples, triples with
source text and plain retrieval using larger models, and \citet{kelesis2026} show that small-model
extraction lowers graph-reasoning accuracy. This chapter adds a repair sequence that changes one
stage of a small-model graph pipeline at a time, on the same 200 questions. It uses that sequence
to locate where evidence is lost: in extraction, in selection, in the triple representation or in
the reader. The individual repairs are not new. Using a graph together with passage retrieval has
precedents (Section~\ref{sec:ch6_scope}), and attaching source text to triples is one of the
conditions compared by \citet{rag_vs_graphrag}. The contribution
is applying them one at a time to a small-model pipeline, so that each loss of evidence can be
assigned to one stage. Two controls complete the sequence. An extraction-budget control shows a floor below
which retrieving more facts does not recover answers, and a reader--extractor crossover separates
the model that builds the graph from the model that answers.

\section{What the graph baseline represents}
\label{sec:ch6_scope}

The graph arm is the SLM-extracted dialogue pipeline defined in Chapter~\ref{ch:methodology},
not an unchanged Microsoft GraphRAG installation. It stores extracted triples as directed
edges, merging repeated relations. Its results depend both on which facts survive
construction and on how evidence reaches the reader. Graph-guided retrieval and presenting
triples as evidence are distinct choices: HippoRAG, for example, uses a graph to rank
source passages~\citep{hipporag2024}. The repairs below test these choices within our pipeline.

\section{The repair sequence}
\label{sec:ch6_ladder}

The versions are called GraphRAG v1 to v4 because each repair changes one stage of the whole
retrieval pipeline (construction, selection or reader input), not the form of the graph. Every
version stores (subject, relation, object) edges in a directed multigraph. Table~\ref{tab:graph_versions} lists what each version changes. Version 2 changes how
triples are selected, and at the same time stores source links and standardizes names. Version 3
raises the extraction budget, which rebuilds the graph with more triples. Version 4 keeps the graph
and the selected triples of version 3 and changes only the text given to the reader.

Version 1 was the original implementation. Each later version addresses the failure that remained
visible in the one before it: selection that depended on exact name matches (v2), cut-off extraction
output (v3) and missing detail in the triples (v4). Source restoration therefore came last. While
selection and extraction still lost evidence, its effect could not have been separated from theirs.

The main comparisons use 200 category-balanced LoCoMo questions under the workstation
protocol of Chapter~\ref{ch:methodology}; earlier 100-question runs guided development.
Version 1 performs no embedding; the other conditions follow the reporting convention
in Section~\ref{sec:retrieval}.

\subsection{GraphRAG v1: entity-match retrieval}

The original implementation extracted triples from batches of up to five dialogue turns. Each
extraction call had a 256-token output limit. At question time, the system searched for graph
nodes whose names matched terms in the question, then collected relationships within one link
of the matching nodes. It ranked those relationships using term matches and simple rules for
dates and session information, and supplied up to 12 triples to the reader.

This baseline could fail during extraction or during selection. Failed extraction removed
facts before retrieval, while name matching and score ties could exclude useful stored
relationships. The following case illustrates both problems.

\subsubsection{A traced retrieval failure}
\label{sec:ch6_retrieval_failure}

Figure~\ref{fig:ch4_extracted_subgraph_case_study} traces ``What did Caroline research?''
on LoCoMo \texttt{conv-26}, using the early phone graph with a 256-token extraction limit
and v1 one-hop retrieval. That graph has 148 nodes and 154 edges; the repaired workstation
\emph{v3} graph of the same conversation has 537 nodes and 551 edges. Phone extraction
failed on 58 of 84 batches, compared with 9 of 1,179 across all ten workstation E2B v3
graphs. Speaker-incident edges account for 85\% of this phone graph and 81\% of this
single workstation v3 graph. The ten-history v3 aggregate is reported separately in
Section~\ref{sec:graph_topology}; the phone pilot is not the graph behind the workstation
accuracy grid.

\begin{figure}[htbp]
  \centering
  \includegraphics[width=\textwidth]{results/thesis_visuals/fig21_extracted_subgraph_case_study.pdf}
  \caption{Retrieval failure with the v1 one-hop rule in the early Gemma~4 E2B phone graph for
  LoCoMo \texttt{conv-26}. (A) Scores
  of 77 one-hop candidate facts; colours identify the 12 facts included in the recorded prompt.
  Seventy-two candidates tie, making selection depend on iteration order. (B) The path from
  omitted evidence to the incorrect answer. Workstation vector retrieval also misses the
  evidence turn, so this case does not demonstrate vector superiority.}
  \label{fig:ch4_extracted_subgraph_case_study}
\end{figure}

The labelled evidence turn says ``Researching adoption agencies'', but no extracted
relation records that fact. Retrieval matched the speaker hub and tied 72 of its 77
candidate facts. The recorded prompt contained nine exchanges between the speakers
among its twelve facts, omitted the tied adoption fact, and produced the wrong answer.
The figure traces each stage; adding the adoption fact would not itself restore the
missing research relation.

The tie order follows Python's per-process hash seed. Repeating the query under 50 seeds
kept the adoption triple in 31 runs and reproduced the pilot's twelve facts in 19.
In the workstation v1 runs, 164 of 200 questions per Gemma reader had ties across the
12-fact cutoff, including 138 shared cases. Their reported v1 accuracies are single draws
from that variation. Vector-first retrieval from v2 onward returned identical contexts
under three tested seeds (\texttt{results/ego\_tie\_order.json}). Workstation vector
retrieval also missed this question's evidence at every tested $k$ from 2 to 50; this
case does not demonstrate that vector retrieval would answer it correctly.

\subsection{GraphRAG v2: turn-first retrieval}

Version 2 began with the same kind of vector search used by the vector-retrieval baseline.
The system embedded each dialogue turn, selected the ten turns most similar to the question,
and considered graph edges whose recorded source batches overlapped those turns. It ranked
candidate edges by the similarity of their best matching turn, multiplied by the edge's stored
weight, and kept up to 12.

This change reduced dependence on matching the question to the exact wording of graph nodes.
It also stored links back to the dialogue and standardized some names, for example by treating
``mom'' and ``mother'' alike. Because these changes were introduced together, version 2 is a
combined retrieval and naming repair. Its accuracy change cannot be assigned to vector search
alone.

The source links have a further limitation. Each edge records the range of turns in the batch
from which it was extracted. It does not record which individual sentence supports that edge.
A selected edge can therefore be associated with a retrieved turn from the same batch even
when that turn expresses a different fact. This detail affects both source restoration in
version 4 and the later SPE design.

Version 2 still showed only triples to the answer model. On the 200 E2B questions, accuracy
rose from 15.0\% to 18.5\%. The extraction counters remained at 507 failed
batches out of 1,179. Better selection could only work with the facts that had entered the graph.

\subsection{GraphRAG v3: larger extraction budget}

The failure logs included structured responses cut off before a triple or the enclosing object
was complete. An output limit of 256 tokens could be too short for the relationships in five
dialogue turns. Version 3 retained version 2's retrieval method and doubled the extraction
limit to 512 tokens.

For E2B, failed batches fell from 507 to 9 out of 1,179, or from 43.0\% to 0.8\%. The number
of extracted triples rose from 3,531 to 7,699. Accuracy moved from 18.5\% to 21.5\%, a modest
difference that did not separate statistically. The clear
effect of version 3 was on graph construction, not on the answer score.

The large reduction in failed batches establishes that the output budget was a practical
problem for this extractor. It does not establish that every original failure had the same
cause. The original version 2 log contains 504 JSON decoding errors, where the response could not
be read as a structured object, and three type errors caused by fields with the wrong value type.
The version 3 log contains four JSON decoding errors and five type errors. These counters
record whether the extraction pipeline could process a response, not whether every resulting
fact was correct or complete.

\subsection{GraphRAG v4: source turns restored}

After version 3, E2B could build a substantially larger graph with few failed batches, yet its
answer accuracy remained below vector retrieval. Version 4 tested whether useful details were
missing from the representation given to the reader. It kept the same graph and selected
triples as version 3, then appended retrieved dialogue turns whose indices overlapped the
selected edges' source batches. Each source turn was included once.

This comparison changes the text available for answering without changing the graph contents or
changing which triples are selected. It can test whether the attached text helps. Because the
source links are recorded by batch, it cannot establish that each added detail was a faithful
expansion of one particular selected triple.

The reported runs use the corrected graph loader; an earlier attempt that failed to apply
source expansion was excluded (Appendix~\ref{app:repro_graph_loader}).

\section{How much the repairs recovered}
\label{sec:ch6_repair_results}

Figure~\ref{fig:ch6_repair_ladder} shows the E2B sequence on 200 questions. Accuracy increased
from 15.0\% in version 1 to 33.0\% in version 4. The largest individual increase, 11.5
percentage points, followed source restoration. Mean prompt length stayed near 520 tokens
through version 3 and rose to 813 tokens in version 4. Vector retrieval scored 36.0\% with a
mean prompt length of 824 tokens.

\begin{figure}[htbp]
  \centering
  \includegraphics[width=\textwidth]{results/thesis_visuals/fig03_graph_repair.pdf}
  \caption{E2B repair sequence on the same 200 LoCoMo questions, using the record runs. Version 1
  retrieves by matching names and embeds nothing; versions 2 to 4 and vector retrieval use the
  corrected task prefixes. The accuracy intervals are Wilson 95\% intervals. Prompt tokens are
  the measured whole-prompt counts on the workstation. Table~\ref{tab:ch6_primary_results}
  extends versions 3 and 4 and vector retrieval to the other four readers.}
  \label{fig:ch6_repair_ladder}
\end{figure}

Source restoration helped four of five readers (Table~\ref{tab:ch6_primary_results}).
For E2B, 27 answers changed from wrong to correct and four from correct to wrong,
giving an 11.5-point gain that survives the designated Holm correction ($p=0.00017$).
The E4B gain is five points (exploratory, raw $p=0.0129$).

\begin{table}[htbp]
  \centering
  \small
  \setlength{\tabcolsep}{5pt}
  \begin{tabular}{lrrrr}
    \toprule
    Model & v3 triples & v4 triples + text & Vector retrieval & v4 gain \\
    \midrule
    Gemma 4 E2B  & 21.5\% & 33.0\% & 36.0\% & +11.5 \\
    Gemma 4 E4B  & 27.5\% & 32.5\% & 40.5\% & +5.0 \\
    Qwen3.5 0.8B &  6.0\% &  9.0\% & 22.0\% & +3.0 \\
    Qwen3.5 2B   &  9.5\% &  9.0\% & 32.5\% & -0.5 \\
    Qwen3.5 4B   & 23.5\% & 32.0\% & 39.0\% & +8.5 \\
    \bottomrule
  \end{tabular}
  \caption{Primary workstation results with corrected embedding prefixes, 200 paired LoCoMo
  questions per condition. Versions 3 and 4 use a 512-token extraction limit and allow up to
  12 selected triples. The final column is the change from v3 to v4 in percentage points.
  Qwen rows are descriptive checks in a second model family.}
  \label{tab:ch6_primary_results}
\end{table}

Qwen 4B gained 8.5 percentage points, whereas Qwen 0.8B remained weak and Qwen 2B was
essentially unchanged. The latter's extraction failed on 84.6\% of batches
(Table~\ref{tab:ch6_extraction}), limiting the evidence available for restoration.

By question category, source restoration mainly helps single-hop questions: it fixes 13
and breaks one for E2B, and fixes seven and breaks none for E4B. These category comparisons
are exploratory.

No source-restored point estimate reached vector retrieval. E2B's remaining three-point
gap did not separate statistically; E4B retained an eight-point gap (exploratory
$p=0.0070$). E2B mean prompt length rose from 519.6 to 812.7 tokens, close to vector
retrieval's 824.4.

\section{What reached the answer model}
\label{sec:ch6_evidence}

The saved prompts allow a more direct inspection than the accuracy totals alone. Consider
the LoCoMo question ``What martial arts has John done?'' The reference answer names kickboxing
and taekwondo. Version 3 supplied 12 triples, none of which named either sport, and E2B replied
``No information available.'' One selected triple referred only to a question about a martial
arts place or yoga studio.

Version 4 supplied the same 12 triples and added five dialogue turns. Two contained John's
statements ``I'm doing kickboxing'' and ``I'm off to do some taekwondo!'' E2B then answered
``kickboxing and taekwondo,'' which the judge accepted. This example shows information absent
from the selected triple text becoming available through source restoration. It also shows why
batch-level source links must be described carefully: useful text can accompany a triple
without expressing that triple's particular relationship. The example illustrates a mechanism;
it does not estimate how often that mechanism explains the full set of errors.

We also checked whether the prompts contained dialogue turns identified as evidence by
LoCoMo. This analysis covered 199 questions with usable evidence labels; one question was
excluded from this diagnostic, while all 200 remained in the accuracy calculation.
Table~\ref{tab:ch6_evidence} reports the proportion of questions for which at least one
labelled evidence turn appeared in the prompt. It does not measure whether all the evidence
needed for an answer was present.

\begin{table}[htbp]
  \centering
  \small
  \begin{tabular}{lrr}
    \toprule
    E2B condition & At least one evidence turn & Reference answer string \\
    \midrule
    v3 triples           & 1/199 (0.5\%)   & 33/199 (16.6\%) \\
    v4 triples + text    & 97/199 (48.7\%) & 40/199 (20.1\%) \\
    Vector retrieval     & 127/199 (63.8\%) & 42/199 (21.1\%) \\
    \bottomrule
  \end{tabular}
  \caption{Prompt checks for the primary E2B conditions. The first diagnostic searches for
  the text of dataset-labelled evidence turns. The second searches the normalized prompt for
  the reference answer string. Both are text-matching checks, not semantic assessments of
  whether the prompt is sufficient to answer.}
  \label{tab:ch6_evidence}
\end{table}

The near-zero evidence-turn match for version 3 is expected because triples reword the
dialogue. It does not mean the graph contains no useful facts. The answer-string check is
also limited: it can miss paraphrases and differently written dates, and it searches the
whole prompt, including the question. We therefore use these measures as supporting
diagnostics. The paired source-restoration experiment and the inspected example provide
stronger evidence that useful detail was absent from the text supplied by version 3.

All five readers' graph and vector conditions start from the same ten retrieved turns.
For E2B, these contain labelled evidence on 127 questions, of which the twelve-triple
selection retains evidence-linked triples on 97. On those 97 questions, v3 answers 26
correctly, while v4 and vector retrieval both answer 44 correctly. Within the 127-question
subset, the remaining gap is in the 30 questions whose evidence the graph selection drops:
v3, v4 and vector retrieval answer 4, 6 and 10 correctly. Separately, on the 72 questions
with labelled evidence that retrieval misses, v4 answers 15 correctly and vector retrieval
17. These group counts locate the observed shortfall; they do not isolate the causes of
errors.

When complete labelled evidence is present, the two methods have similar conditional
accuracy, without a consistent advantage across readers (Table~\ref{tab:ch4_complete_use}).
The groups contain different questions, so the comparison is descriptive.

\begin{table}[htbp]
  \centering
  \small
  \setlength{\tabcolsep}{7pt}
  \begin{tabular}{lrrrr}
    \toprule
    & \multicolumn{2}{c}{Graph + sources} & \multicolumn{2}{c}{Vector retrieval} \\
    \cmidrule(lr){2-3} \cmidrule(lr){4-5}
    Reader & Complete cases & Accuracy & Complete cases & Accuracy \\
    \midrule
    Gemma E2B & 67 & 56.7\% & 88 & 50.0\% \\
    Gemma E4B & 67 & 53.7\% & 88 & 56.8\% \\
    \bottomrule
  \end{tabular}
  \caption{Accuracy when the complete labelled evidence is present. The case counts differ because
           graph source restoration and vector retrieval do not cover the same number of questions.}
  \label{tab:ch4_complete_use}
\end{table}

\section{Would more triples be enough?}
\label{sec:ch6_depth}

The default limit of 12 triples was another possible constraint. We therefore reused the
version 3 graphs and varied the maximum number of triples placed in the answer prompt.
The main extraction settings stayed fixed. Table~\ref{tab:ch6_depth} shows the results with
corrected embedding prefixes.

\begin{table}[htbp]
  \centering
  \scriptsize
  \setlength{\tabcolsep}{4pt}
  \resizebox{\textwidth}{!}{%
  \begin{tabular}{lrrrrrrrr}
    \toprule
    & \multicolumn{2}{c}{4 triples} & \multicolumn{2}{c}{12 triples}
    & \multicolumn{2}{c}{24 triples} & \multicolumn{2}{c}{Vector retrieval} \\
    \cmidrule(lr){2-3}\cmidrule(lr){4-5}\cmidrule(lr){6-7}\cmidrule(lr){8-9}
    Reader & Accuracy & Tokens & Accuracy & Tokens & Accuracy & Tokens & Accuracy & Tokens \\
    \midrule
    Gemma E2B & 13.0\% & 287 & 21.5\% & 520 & 29.5\% & 870 & 36.0\% & 824 \\
    Gemma E4B & 16.0\% & 295 & 27.5\% & 539 & 29.5\% & 909 & 40.5\% & 824 \\
    Qwen 4B & 11.5\% & 283 & 23.5\% & 513 & 29.0\% & 864 & 39.0\% & 812 \\
    Qwen 0.8B & n/a & n/a & 6.0\% & --- & 7.0\% & --- & 22.0\% & --- \\
    Qwen 2B & n/a & n/a & 9.5\% & --- & 10.0\% & --- & 32.5\% & --- \\
    \bottomrule
  \end{tabular}%
  }
  \caption{Graph-depth controls on 200 paired LoCoMo questions. Each reader's v3 graph
  is fixed; only the maximum number of retrieved triples changes. Tokens are mean
  whole-prompt counts. Vector retrieval is the reference from Chapter~\ref{ch:memory_evaluation}.
  The four-triple condition was not run for the two smallest Qwen models (n/a);
  dashes indicate token means omitted from this table.}
  \label{tab:ch6_depth}
\end{table}

Increasing the limit from twelve to twenty-four triples gives E2B 18 improvements and two
regressions, an eight-point gain (exploratory $p=0.00040$). Its mean prompt grows from
519.6 to 870.2 tokens, yet its 29.5\% accuracy remains below source restoration's 33.0\%
at 812.7 tokens. This is a descriptive comparison between different evidence sets.
E4B gains two points and Qwen 4B gains 5.5; the smaller Qwen models change by only one
or two answers (raw $p=0.5$ for 0.8B and $p=1$ for 2B).

Reducing the limit from twelve to four costs E2B 23 correct answers while gaining six
(raw $p=0.002316$). E4B loses 28 and gains five; Qwen 4B loses 26 and gains two.
With losses at four and gains at twenty-four for E2B and Qwen 4B, twelve triples does
not saturate observed accuracy for those readers. None of the tested limits
for these three readers reaches the accuracy of source restoration.

The answer-cap and reasoning controls show no consistent recovery from additional output
tokens under the tested settings (Section~\ref{sec:ch4_reasoning}).

\section{Extraction quality and model choice}
\label{sec:ch6_extractors}

The Gemma repair did not establish that a 512-token limit would suit every model.
Table~\ref{tab:ch6_extraction} compares extraction failures and output volume across the
five models, including the 1,024-token controls for Qwen 0.8B and 2B. All conditions
processed the same 1,179 batches from the ten LoCoMo conversations.

\begin{table}[htbp]
  \centering
  \small
  \setlength{\tabcolsep}{5pt}
  \begin{tabular}{lrrrrr}
    \toprule
    Extractor & Output limit & Failed batches & Extracted triples & Stored edges & QA accuracy \\
    \midrule
    Gemma 4 E2B  & 512   & 9 (0.8\%)    & 7,699 & 7,458 & 21.5\% \\
    Gemma 4 E4B  & 512   & 7 (0.6\%)    & 6,264 & 6,228 & 27.5\% \\
    Qwen3.5 4B   & 512   & 45 (3.8\%)   & 7,699 & 7,292 & 23.5\% \\
    Qwen3.5 0.8B & 512   & 931 (79.0\%) & 1,898 & 1,678 & 6.0\% \\
    Qwen3.5 0.8B & 1,024 & 889 (75.4\%) & 2,742 & 2,427 & 6.5\% \\
    Qwen3.5 2B   & 512   & 998 (84.6\%) & 1,925 & 1,853 & 9.5\% \\
    Qwen3.5 2B   & 1,024 & 802 (68.0\%) & 5,683 & 5,323 & 11.0\% \\
    \bottomrule
  \end{tabular}
  \caption{Extraction-budget controls over 1,179 LoCoMo batches. Output limits are tokens per call.
  Failed batches cannot be processed; valid empty returns are separate. Triples are counted
  before repeated edges merge. E2B and Qwen 4B happen to produce the same number of triples;
  their stored edges differ. QA uses v3 retrieval, corrected prefixes, and the same model for
  extraction and answering. No transport errors were recorded in the inspected logs.}
  \label{tab:ch6_extraction}
\end{table}

Doubling the Qwen 2B extraction allowance nearly tripled its output, from 1,925 to 5,683
triples. The number of stored graph edges rose from 1,853 to 5,323. Accuracy changed from
9.5\% to 11.0\%, a net gain of three answers, which did not separate statistically and is
too small to interpret as a recovery. Qwen 0.8B rose only from 6.0\% to 6.5\%
with the larger allowance. Neither result reproduced the near-elimination of
failed batches seen with Gemma E2B.

These controls show that extraction volume alone was insufficient to recover answer
quality in the tested small Qwen configurations. They do not prove that extraction quality
is unimportant. Many batches still failed at 1,024 tokens, and we did not measure the
correctness or completeness of every surviving triple. Within Qwen, the 4B condition was
the first tested size with a low failure rate at the 512-token setting. This is a boundary
observed for this prompt, budget, model family, and runtime, rather than a universal model-size
threshold for graph construction.

\subsection{Structure of the repaired Gemma graphs}
\label{sec:graph_topology}

Table~\ref{tab:graph_topology} describes the workstation E2B and E4B \emph{v3} caches
across ten LoCoMo conversations. These repaired graphs also supply the SPE evaluation.

\begin{table}[htbp]
  \centering
  \scriptsize
  \setlength{\tabcolsep}{3pt}
  \begin{tabular}{lrrrrrcc}
    \toprule
    Extractor & Nodes ($|V|$) & Edges ($|E|$) & Degree ($\bar{k}$) & Components & Giant \% & Empty \% & Extraction failure \% \\
    \midrule
    E2B & $676.6 \pm 140.1$ & $745.8 \pm 151.0$ & 2.21 & $36.5 \pm 12.2$ & 86.1\% & 0.00\% & 0.76\% \\
    E4B & $588.2 \pm 108.0$ & $622.8 \pm 115.2$ & 2.12 & $24.3 \pm 8.6$  & 89.7\% & 7.12\% & 0.59\% \\
    \bottomrule
  \end{tabular}
  \caption{Properties of the repaired Gemma~4 v3 graphs used in the SPE evaluation across 10 LoCoMo
           conversations. Nodes, edges and components are per-conversation means $\pm$ standard
           deviations. Degree is $2|E|/|V|$ counting every stored edge; read as a simple undirected
           graph, the mean degree is 1.96 (E2B) and 1.97 (E4B). Empty \% and Extraction failure \% are
           pooled over all 1,179 extraction batches.}
  \label{tab:graph_topology}
\end{table}

Across the ten repaired workstation E2B v3 graphs, the mean share of edges touching
either speaker is 83\%.
For each question, \texttt{graph\_audit.py} takes the entities of all triples whose source
batches overlap the ten retrieved turns, then counts graph edges incident to those entities
(one step). This edge share averages 87\% across questions. Excluding speakers as stepping
stones, the first step adds under four percentage points to the seed triples' edge share
and the second adds less than one. These observations describe graph reachability;
a traversal policy designed for this structure was not tested.

E4B has fewer disconnected components and a larger giant-component share than E2B.
Only 1.7\% of E2B edges are reinforced across multiple dialogue turns, so repeated
relations contribute little consolidation in this representation. These measurements
describe the graph inputs; they do not establish why one extractor answers better.

\subsection{Empty responses and discarded values}
\label{sec:ch6_extraction_replay}

The original E4B v3 extraction returned 84 empty batches out of 1,179 (7.12\%), against
none for E2B. Because the original counter did not retain those responses, all 1,179 E4B
calls were replayed with instrumentation. Total extraction volume reproduced within 0.6\%;
all 96 replayed empty batches were explicit \texttt{\{"triples": []\}} returns.
The count difference is treated as runtime variation.

Replacing each original empty batch with E4B's mean non-empty yield would recover about
484 of the 1,435-triple difference from E2B, roughly one third. E4B also produces fewer
triples in non-empty batches. Inspection found both appropriate omissions of social
exchanges and missed facts such as named pets, purchases and dated activities. Empty
returns therefore do not establish better precision, nor does this audit estimate
extraction recall or explain downstream accuracy.

The replay also found that a rule removing list numbers discarded three bare-numeric
triples out of 6,231 returned triples. It caused none of the empty batches. The stored
graphs were not rebuilt with a revised normalizer, so this diagnostic bounds the defect's
observed effect in the replay rather than establishing that all implementation errors
were removed.

\subsection{Separate the model that builds the graph from the model that answers}

E4B extracted fewer triples than E2B but achieved higher version 3 accuracy. In the usual
runs, however, the same model both built the graph and answered the questions. We separated
those roles by giving each reader the graph produced by each extractor
(Table~\ref{tab:ch6_crossover}). The retrieval settings and question set stayed fixed.

\begin{table}[htbp]
  \centering
  \small
  \begin{tabular}{lrr}
    \toprule
    Answer model & Graph built by E2B & Graph built by E4B \\
    \midrule
    Gemma 4 E2B & 21.5\% & 26.0\% \\
    Gemma 4 E4B & 22.0\% & 27.5\% \\
    \bottomrule
  \end{tabular}
  \caption{Reader-extractor crossover on the same 200 LoCoMo questions, using version 3
  and corrected embedding prefixes. Rows hold the reader fixed; columns hold the graph
  producer fixed. With the graph fixed, both readers received the same graph cache and
  retrieved context, and retrieval metadata matched for all 200 questions.}
  \label{tab:ch6_crossover}
\end{table}

Changing the extractor from E2B to E4B gains 4.5--5.5 points with the reader fixed;
changing only the reader gains 0.5--1.5 points. Most of the difference in these four
cells is therefore associated with the graph producer. The crossover does not identify
whether fact selection, naming or edge ranking is responsible, and it does not show
that reader size is generally unimportant. Graph size and empty-batch frequency alone
cannot explain the advantage.

\section{What the study establishes and where it stops}
\label{sec:ch6_limits}

The experiments distinguish failed or omitted extraction, evidence excluded during
selection, detail missing from selected triples, and errors after evidence reaches the
reader. These constraints should not be collapsed into one explanation. Larger-model
work likewise finds that graph retrieval design and the supplied representation both
affect outcomes~\citep{rag_vs_graphrag}.

The depth controls are not exactly token-matched. Source links identify batches,
normalization defects remain, and there is no complete manually verified reference graph.
The results therefore do not assign all remaining error to triple compression.

Repairs were developed while inspecting this benchmark, so the repair sequence is a
development record on LoCoMo, not a test held out from it; Section~\ref{sec:ch4_lme} repeats
the repaired pipeline on LongMemEval-S. The repaired graph was not
built on the phone; the available phone pilot is reported in
Chapter~\ref{ch:on_device_memory}.

\section{From full source restoration to selective expansion}
\label{sec:ch6_to_spe}

Source restoration improves answers but increases prompt length. This motivates choosing
whether each retrieved fact should be omitted, retained as a triple or accompanied by
source text under a shared context budget. Chapter~\ref{ch:spe} develops Selective
Provenance Expansion and evaluates that decision against criteria fixed before its runs.
